% Examen ALEA du 9 janvier 2026, exercice 2.
\uuid{fyMW}
\titre{Densité conjointe sur le carré unité}
\niveau{L2}
\module{Probabilité}
\chapitre{Probabilité continue}
\sousChapitre{Couples de variables aléatoires}
\theme{Densité conjointe, lois marginales, indépendance}
\auteur{}
\datecreate{2026-10-01}
\organisation{AMSCC}
\difficulte{}
\contenu{
\texte{Soit $a\in\mathbb{R}$ et $(X,Y)$ un couple de variables aléatoires admettant pour densité
\[f(x,y)=a(x^2+y^2)\mathbf{1}_{[0,1]}(x)\mathbf{1}_{[0,1]}(y).\]}
\begin{enumerate}
  \item \question{Déterminer $a$.}
  \item \question{Déterminer les lois marginales du couple $(X,Y)$.}
  \item \question{Calculer les espérances $\mathbb{E}[X]$, $\mathbb{E}[Y]$ et les variances $\sigma^2(X)$, $\sigma^2(Y)$.}
  \item \question{Les variables $X$ et $Y$ sont-elles indépendantes ?}
\end{enumerate}
}
